\documentclass[12pt]{article}
\usepackage[a4paper,margin=1in]{geometry}
\usepackage[utf8]{inputenc}
\usepackage[T1]{fontenc}
\usepackage{amsmath,amssymb}
\usepackage{textcomp}
\usepackage{parskip}
\begin{document}
\begin{center}
{\Large \textbf{2020 STEP II}}
\end{center}

\subsection*{Page 1}
\begin{flushleft}
\bigskip\textbf{Section A: Pure Mathematics}\par
\bigskip\textbf{Question 1.} (i) Use the substitution x = , where 0 < u < 1, to find in terms of x the integral\par
1 - u\par
\[\int\]
\[dx (where x > 1).\]
3 1\par
x 2 (x - 1) 2\par
\textbf{(ii)} Find in terms of x the integral\par
\[\int\]
\[dx (where x > 2).\]
3 1\par
(x - 2) 2 (x + 1) 2\par
\textbf{(iii)} Show that\par
Z \infty 1\par
\[dx = 1 \pi.\]
1 1 3\par
2 (x - 1)(x - 2) 2 (3x - 2) 2\par
\bigskip\textbf{Question 2.} The curves C and C both satisfy the differential equation\par
1 2\par
dy kxy - y\par
\[= ,\]
dx x - kxy\par
\[where k = ln 2.\]
All points on C have positive x and y co-ordinates and C passes through (1 , 1). All points\par
1 1\par
on C have negative x and y co-ordinates and C passes through (-1 , -1).\par
2 2\par
\textbf{(i)} Show that the equation of C can be written as (x - y)2 = (x + y)2 - 2x+y.\par
Determine a similar result for curve C .\par
\[Hence show that y = x is a line of symmetry of each curve.\]
\textbf{(ii)} Sketch on the same axes the curves y = x2 and y = 2x, for x > 0. Hence show that C\par
\[lies between the lines x + y = 2 and x + y = 4.\]
Sketch curve C .\par
\textbf{(iii)} Sketch curve C .\par
\end{flushleft}
\subsection*{Page 2}
\begin{flushleft}
\bigskip\textbf{Question 3.} A sequence u , u , . . . , u of positive real numbers is said to be unimodal if there is a value\par
1 2 n\par
k such that\par
u 6 u 6 . . . 6 u\par
1 2 k\par
and\par
\[u > u > . . . > u .\]
k k+1 n\par
So the sequences 1, 2, 3, 2, 1; 1, 2, 3, 4, 5; 1, 1, 3, 3, 2 and 2, 2, 2, 2, 2 are all unimodal, but\par
1, 2, 1, 3, 1 is not.\par
A sequence u , u , . . . , u of positive real numbers is said to have property L if u u 6 u2\par
1 2 n r-1 r+1 r\par
for all r with 2 6 r 6 n - 1.\par
\textbf{(i)} Show that, in any sequence of positive real numbers with property L,\par
\[u > u =⇒ u > u .\]
r-1 r r r+1\par
Prove that any sequence of positive real numbers with property L is unimodal.\par
\textbf{(ii)} A sequence u , u , . . . , u of real numbers satisfies u = 2αu - α2u for 3 6 r 6 n,\par
1 2 n r r-1 r-2\par
where α is a positive real constant. Prove that, for 2 6 r 6 n,\par
\[u - αu = αr-2(u - αu )\]
r r-1 2 1\par
and, for 2 6 r 6 n - 1,\par
\[u^2 - u u = (u - αu )2.\]
r r-1 r+1 r r-1\par
Hence show that the sequence consists of positive terms and is unimodal, provided\par
\[u > αu > 0.\]
2 1\par
\[In the case u = 1 and u = 2, prove by induction that u = (2 - r)αr-1 + 2(r - 1)αr-2.\]
1 2 r\par
\[Let α = 1 - , where N is an integer with 2 6 N 6 n.\]
N\par
\[In the case u = 1 and u = 2, prove that u is largest when r = N .\]
1 2 r\par
\end{flushleft}
\subsection*{Page 3}
\begin{flushleft}
\bigskip\textbf{Question 4.} (i) Given that a, b and c are the lengths of the sides of a triangle, explain why c < a + b,\par
\[a < b + c and b < a + c.\]
\textbf{(ii)} Use a diagram to show that the converse of the result in part (i) also holds: if a, b and\par
\[c are positive numbers such that c < a + b, a < b + c and b < c + a then it is possible\]
to construct a triangle with sides of length a, b and c.\par
\textbf{(iii)} When a, b and c are the lengths of the sides of a triangle, determine in each case\par
whether the following sets of three lengths can\par
• always\par
• sometimes but not always\par
• never\par
form the sides of a triangle. Prove your claims.\par
(A) a + 1, b + 1, c + 1.\par
a b c\par
(B) , , .\par
b c a\par
(C) |a - b|, |b - c|, |c - a|.\par
\[(D) a^2 + bc, b^2 + ca, c^2 + ab.\]
\textbf{(iv)} Let f be a function defined on the positive real numbers and such that, whenever\par
\[x > y > 0,\]
f(x) f(y)\par
\[f(x) > f(y) > 0 but < .\]
x y\par
Show that, whenever a, b and c are the lengths of the sides of a triangle, then f(a), f(b)\par
and f(c) can also be the lengths of the sides of a triangle.\par
\end{flushleft}
\subsection*{Page 4}
\begin{flushleft}
5 If x is a positive integer, the value of the function d(x) is the sum of the digits of x in base\par
\[10. For example, d(249) = 2 + 4 + 9 = 15.\]
n-1\par
X\par
\[An n-digit positive integer x is written in the form a \times 10r, where 0 6 a 6 9 for all\]
r r\par
\[r=0\]
\[0 6 r 6 n - 1 and a > 0.\]
n-1\par
\textbf{(i)} Prove that x - d(x) is non-negative and divisible by 9.\par
\textbf{(ii)} Prove that x - 44d(x) is a multiple of 9 if and only if x is a multiple of 9.\par
\[Suppose that x = 44d(x). Show that if x has n digits, then x 6 396n and x > 10n-1,\]
and hence that n 6 4.\par
\[Find a value of x for which x = 44d(x). Show that there are no further values of x\]
satisfying this equation.\par
\textbf{(iii)} Find a value of x for which x = 107d (d(x)). Show that there are no further values of\par
x satisfying this equation.\par
\lambda \mu\par
a b\par
\bigskip\textbf{Question 6.} A 2 \times 2 matrix M is real if it can be written as M = , where a, b, c and d are real.\par
c d\par
\[In this case, the trace of matrix M is defined to be tr(M) = a + d and det(M) is the\]
\[determinant of matrix M. In this question, M is a real 2 \times 2 matrix.\]
\textbf{(i)} Prove that\par
\[tr(M^2) = tr(M)2 - 2det(M).\]
\textbf{(ii)} Prove that\par
\[M^2 = I but M 6= ±I ⇐⇒ tr(M) = 0 and det(M) = -1,\]
and that\par
\[M^2 = -I ⇐⇒ tr(M) = 0 and det(M) = 1.\]
\textbf{(iii)} Use part (ii) to prove that\par
\[M^4 = I ⇐⇒ M^2 = ±I.\]
\[Find a necessary and sufficient condition on det(M) and tr(M) so that M^4 = -I.\]
\textbf{(iv)} Give an example of a matrix M for which M8 = I, but which does not represent a\par
rotation or reflection. [Note that the matrices ±I are both rotations.]\par
\end{flushleft}
\subsection*{Page 5}
\begin{flushleft}
\bigskip\textbf{Question 7.} In this question, w = .\par
z - 2\par
\textbf{(i)} Let z be the complex number 3 + ti, where t \in R. Show that |w - 1| is independent\par
\[of t. Hence show that, if z is a complex number on the line Re(z) = 3 in the Argand\]
diagram, then w lies on a circle in the Argand diagram with centre 1.\par
\[Let V be the line Re(z) = p, where p is a real constant not equal to 2. Show that, if z\]
lies on V , then w lies on a circle whose centre and radius you should give in terms of\par
\[p. For which z on V is Im(w) > 0?\]
\textbf{(ii)} Let H be the line Im(z) = q, where q is a non-zero real constant. Show that, if z lies\par
on H, then w lies on a circle whose centre and radius you should give in terms of q.\par
\[For which z on H is Re(w) > 0?\]
\bigskip\textbf{Question 8.} In this question, f(x) is a quartic polynomial where the coefficient of x4 is equal to 1, and\par
\[which has four real roots, 0, a, b and c, where 0 < a < b < c.\]
Z x\par
\[F(x) is defined by F(x) = f(t) dt.\]
\[The area enclosed by the curve y = f(x) and the x-axis between 0 and a is equal to that\]
between b and c, and half that between a and b.\par
\textbf{(i)} Sketch the curve y = F(x), showing the x co-ordinates of its turning points.\par
\[Explain why F(x) must have the form F(x) = 1 x^2(x - c)2(x - h), where 0 < h < c.\]
Find, in factorised form, an expression for F(x) + F(c - x) in terms of c, h and x.\par
\textbf{(ii)} If 0 6 x 6 c, explain why F(b) + F(x) > 0 and why F(b) + F(x) > 0 if x 6= a.\par
\[Hence show that c - b = a or c > 2h.\]
\[By considering also F(a) + F(x), show that c = a + b and that c = 2h.\]
\textbf{(iii)} Find an expression for f(x) in terms of c and x only.\par
\[Show that the points of inflection on y = f(x) lie on the x-axis.\]
\end{flushleft}
\subsection*{Page 6}
\begin{flushleft}
\bigskip\textbf{Section B: Mechanics}\par
9 Point A is a distance h above ground level and point N is directly below A at ground level.\par
Point B is also at ground level, a distance d horizontally from N . The angle of elevation of\par
A from B is β. A particle is projected horizontally from A, with initial speed V . A second\par
particle is projected from B with speed U at an acute angle θ above the horizontal. The\par
horizontal components of the velocities of the two particles are in opposite directions. The\par
two particles are projected simultaneously, in the vertical plane through A, N and B.\par
Given that the two particles collide, show that\par
V h\par
\[d sin θ - h cos θ =\]
U\par
and also that\par
\textbf{(i)} θ > β;\par
r\par
gh\par
\textbf{(ii)} U sin θ > ;\par
U\par
\textbf{(iii)} > sin β.\par
V\par
Show that the particles collide at a height greater than 1 h if and only if the particle projected\par
from B is moving upwards at the time of collision.\par
\end{flushleft}
\subsection*{Page 7}
\begin{flushleft}
\bigskip\textbf{Question 10.} A particle P of mass m moves freely and without friction on a wire circle of radius a, whose\par
axis is horizontal. The highest point of the circle is H, the lowest point of the circle is L and\par
\[angle P HL = θ. A light spring of modulus of elasticity λ is attached to P and to H. The\]
natural length of the spring is l, which is less than the diameter of the circle.\par
\textbf{(i)} Show that, if there is an equilibrium position of the particle at θ = α, where α > 0,\par
λl\par
\[then cos α = .\]
2(aλ - mgl)\par
2mgl\par
\[Show also that there will only be such an equilibrium position if λ > .\]
2a - l\par
When the particle is at the lowest point L of the circular wire, it has speed u.\par
\textbf{(ii)} Show that, if the particle comes to rest before reaching H, it does so when θ = β,\par
where cos β satisfies\par
\[mu2\]
\[(cos α - cos β)2 = (1 - cos α)2 + cos α,\]
2aλ\par
λl\par
\[where cos α = .\]
2(aλ - mgl)\par
2aλ\par
\[Show also that this will only occur if u^2 < (2 - sec α).\]
m\par
\end{flushleft}
\subsection*{Page 8}
\begin{flushleft}
\bigskip\textbf{Section C: Probability and Statistics}\par
\bigskip\textbf{Question 11.} A coin is tossed repeatedly. The probability that a head appears is p and the probability\par
\[that a tail appears is q = 1 - p.\]
\textbf{(i)} A and B play a game. The game ends if two successive heads appear, in which case A\par
wins, or if two successive tails appear, in which case B wins.\par
Show that the probability that the game never ends is 0.\par
p\par
Given that the first toss is a head, show that the probability that A wins is .\par
1 - pq\par
Find and simplify an expression for the probability that A wins.\par
\textbf{(ii)} A and B play another game. The game ends if three successive heads appear, in which\par
case A wins, or if three successive tails appear, in which case B wins.\par
Show that\par
\[P(A wins | the first toss is a head) = p^2 + (q + pq) P(A wins | the first toss is a tail)\]
and give a similar result for P(A wins | the first toss is a tail).\par
Show that\par
\[p^2(1 - q^3)\]
\[P(A wins) = .\]
\[1 - (1 - p^2)(1 - q^2)\]
\textbf{(iii)} A and B play a third game. The game ends if a successive heads appear, in which case\par
A wins, or if b successive tails appear, in which case B wins, where a and b are integers\par
greater than 1.\par
Find the probability that A wins this game.\par
\[Verify that your result agrees with part (i) when a = b = 2.\]
\end{flushleft}
\subsection*{Page 9}
\begin{flushleft}
Paper II, 2020 final version 10\par
\bigskip\textbf{Question 12.} The score shown on a biased n-sided die is represented by the random variable X which has\par
\[distribution P(X = i) = + ε for i = 1, 2, . . . , n, where not all the ε are equal to 0.\]
i i\par
n\par
\textbf{(i)} Find the probability that, when the die is rolled twice, the same score is shown on\par
both rolls. Hence determine whether it is more likely for a fair die or a biased die to\par
show the same score on two successive rolls.\par
\textbf{(ii)} Use part (i) to prove that, for any set of n positive numbers x (i = 1, 2, . . . , n),\par
i\par
n i-1 n !2\par
X X n - 1 X\par
x x 6 x .\par
i j i\par
2n\par
\[i=2 j=1 i=1\]
\textbf{(iii)} Determine, with justification, whether it is more likely for a fair die or a biased die to\par
show the same score on three successive rolls.\par
\end{flushleft}
\end{document}
